\section{Consequences for Aharonov--Bohm, Bell, EPR}
\label{sec:consequences_bell_ab_epr_en}

The three preceding structural results --- the holonomy angle as
primitive (Proposition~\ref{prop:theta_p_unique_en}), geometric
decoupling (Theorem~\ref{thm:invariance_geom_en}), and the
arithmetic closed form of the correlations
(Corollary~\ref{cor:closed_form_en}) --- modify the reading of
three classical experiments of foundational physics: the
Aharonov--Bohm configuration, the Bell tests, and the original
EPR configuration. We examine each briefly.

\subsection{Aharonov--Bohm: the phase is $\theta_3$, not $A_\mu$}
\label{ssec:ab_consequence_en}

In the PT reading, the interferometric phase shift
$\Delta\varphi$ observed in an Aharonov--Bohm
configuration~\cite{AharonovBohm1959} is not
$(e/\hbar) \oint A \cdot d\ell$ but $\theta_3$, the holonomy
angle of the modular channel $p = 3$ associated with the flux
enclosed by the particle trajectory. The identification between
the two formulas is tautological in the presence of the DIC
dictionary that translates PT units into SI units; what
distinguishes the PT reading is not the numerical value of the
phase shift but the identification of the \emph{physical support}
of the phase shift.

This identification has two consequences. First, the AB phase
is a strictly topological invariant of the path $\gamma$: it
depends only on the homotopy class of $\gamma$ in the complement
of the flux region, and is insensitive to any continuous
deformation of $\gamma$ within that class. This topology is not
a property of the metric $G(\mu)$; it is anterior to the metric,
in the sense of Theorem~\ref{thm:invariance_geom_en}. Second,
the absence of field $F = 0$ in the interior of an ideal solenoid
is no longer a paradox: $F_{\mu\nu}$ is not the object that
carries the phase information. It is $\theta_3$ that carries it,
and $\theta_3$ is non-zero even where $F$ vanishes, because
$\theta_3$ is defined by the modular connection of PT, which is
not itself a local function of the field.

\subsection{EPR/Bell: intra-channel correlations, not inter-channel}
\label{ssec:bell_consequence_en}

The Bell tests~\cite{Bell1964,Aspect1982,Hensen2015} establish
that the correlations observed between spatially separated
detectors exceed the CHSH bound~\cite{CHSH1969} expected of a
local hidden-variable theory. In the standard reading these
correlations are interpreted as the signature of a fundamental
non-locality of the quantum field: the collapse of an EPR state
upon left-side measurement instantaneously determines the
correlated right-side state, at any distance, and apparently
faster than light (but without allowing superluminal
communication, in accordance with no-signalling).

In the PT reading, this interpretation reverses. The two
measurements performed in a Bell test do \emph{not} act on two
distinct modular channels (in which case, by
Theorem~\ref{thm:invariance_geom_en}, their MI would be
geometry-invariant and of magnitude predicted by the closed
form~\eqref{cor:closed_form_en}). They act on the \emph{same}
modular channel --- typically $p = 2$ for polarisation/spin
experiments, or $p = 3$ for quark-colour experiments
(analogues). An intra-channel correlation is not the object of
Theorem~\ref{thm:invariance_geom_en}: it is a higher-order
correlation in the CRT structure, independent of any notion of
spatial separation, because the modular channel exists before
the construction of space by projection.

This reading is consistent with the empirical signature of the
most precise tests: the correlation observed between two
detectors separated by several kilometres~\cite{Hensen2015} is
strictly identical to that observed at a few metres
\cite{Aspect1982}, to within measurement precision. No significant
dependence on distance has ever been detected. This absence of
dependence is, in the PT reading, the structural expectation:
what is measured is not in space, so spatial separation cannot
modify the result.

\subsection{Dissolution of spatial non-locality}
\label{ssec:dissolution_consequence_en}

The ontological consequence of the two preceding subsections is
the \emph{dissolution} ([DISS] status) of strong-sense spatial
non-locality. The ``paradox'' of information apparently
transmitted faster than $c$ disappears because the correlations
do not propagate in space: they exist in the CRT structure, and
space is merely an emergent projection of that structure via the
metric $G(\mu)$.

The [DISS] status is, in the PT nomenclature, weaker than [THM]
but stronger than mere reinterpretation. A dissolution asserts
that the problem, as posed in the classical framework, is
ill-posed: its implicit premisses --- the existence of a
classical background space and the requirement that every
observed correlation propagate within it --- are not valid in
PT. The problem does not resolve in the sense that one would
have found the ``true'' classical explanation of the EPR
correlations; it dissolves in the sense that the correlations
are seen as objects that never belonged to the spatial domain.

The experimental horizon is then clear: PT predicts no signal
that could separate two spatially separated detectors beyond the
arithmetic CRT correlation, and categorically excludes any
signal that would depend on the metric $G(\mu)$ between the two
detectors. Any positive observation of such a dependence
falsifies PT and the reading presented here simultaneously.
Conversely, experimental confirmation of the geometric invariance
of EPR correlations --- i.e.\ prediction QH1a applied to the
relativistic regime --- would be a direct validation of the PT
reading of the foundational debate.
